\section{Isaac Abarbanel and the Explicit Daylight-Stars Dilemma}

Nearly fifteen hundred years after the close of the Old Testament, the Iberian statesman and biblical commentator Don Isaac Abarbanel (1437--1508) formulated the exact waking contradiction between the stars of Genesis 15:5 and the sunset of Genesis 15:12 in plain, unavoidable Hebrew. Opening his commentary on Genesis 15---a chapter he calls ``great in value and rank'' (\latin{gedolat ha-'erekh we-ha-ma'alah})---Abarbanel announces that after pondering the passage for many days he has drawn up twenty-one prefatory questions (\latin{she'elot}) on its difficulties.\endnote{Isaac Abarbanel, \emph{Perush 'al ha-Torah} on Bereshit 15, opening formula and prefatory Questions 7 and 16 (Warsaw, 1862 printing). Translations made directly from the printed Hebrew text; Abarbanel's citation of \emph{Shabbat} 157 corresponds to modern foliation \emph{Shabbat} 156a.}

\subsection{Question 7: Prophetic Vision, Astrology, and the Waking Dilemma}

Abarbanel devotes Question~7 directly to the opening words of verse 5: ``And he brought him forth outside'' (\latin{wayy\={o}\d{s}\={e}' '\={o}t\={o} ha-\d{h}\={u}\d{s}\={a}h}). Before touching the hour of day, he lays out the two classical Jewish lines of interpretation inherited from earlier centuries---and challenges the visionary reading first:

\begin{itemize}
  \item \textbf{Against the purely visionary reading.} Earlier Jewish commentators (\latin{ke-da'at ha-mefarshim}, ``according to the view of the commentators'') held that the entire episode occurred ``in prophetic vision'' (\latin{be-mar'eh ha-nevu'ah}). Abarbanel asks why Scripture would need to tell us that Abram seemed in his vision to be led out of a tent when the chapter has never mentioned a tent at all; simply saying ``Look now toward heaven'' would have sufficed, since prophecy could come to him in the open field as readily as indoors.
  \item \textbf{The Talmudic departure from astrology.} Perhaps, Abarbanel immediately remarks, that very oddity moved the Talmudic Sages (\emph{Shabbat} 157 in his foliation, corresponding to modern \emph{Shabbat} 156a) to read ``outside'' not as stepping past a tent flap but as leaving astrological determinism behind: \latin{ho\d{s}\={\i}'\={o} m\={e}-'i\d{s}\d{t}agn\={\i}n\={u}t shel\={o}}, ``He brought him out from his astrology.''
\end{itemize}

Having weighed the visionary commentators and the rabbinic reading of \emph{Shabbat} 156a, Abarbanel turns directly to the waking literal alternative and springs a two-horned trap:

\begin{quote}
And if we say that this happened while awake, would that I knew whether it was by day or by night! For it is impossible that it was by day, on account of what is said, ``and count the stars,'' since they are not visible by day; and if it was by night, then what is said [in verse 12] poses a difficulty: ``And it came to pass that the sun was about to set.'' (Abarbanel, Commentary on Genesis 15, Question 7)
\end{quote}

Every element of the modern daylight-stars puzzle stands on the page here in late-fifteenth-century Hebrew---except that Abarbanel draws the exact opposite conclusion from modern daylight-stars preachers. For Abarbanel, a waking daytime command to count the stars is not a hidden test of faith in invisible luminaries; he rejects it as flatly impossible (\latin{'i 'efshar}), ``since they are not visible by day'' (\latin{ki 'einam nir'im ba-yyom}). The chronological difficulty is genuine precisely because the starry sky demands night while verse 12 names the setting sun.

\subsection{Question 16: Sunset, Darkness, and ``a Lamp at Midday''}

In Question~16, Abarbanel presses a companion chronological puzzle on Genesis 15:12 (\latin{wayeh\={\i} ha-shemesh l\={a}v\={o}' we-tard\={e}m\={a}h n\={a}fel\={a}h}, ``And it came to pass, the sun was about to set, and a deep sleep fell''). Nowhere else across all the prophecies of Abraham, he observes, does Scripture specify the hour and time of day when the prophetic influx came upon him; such a time-restriction appears wholly unnecessary.

Abarbanel then cites Rabbenu Nissim of Gerona (the Ran, fourteenth century), who had already noticed this anomaly and proposed an optical reason: God specified the hour because He intended the fiery torch to pass between the divided sacrifices after full nightfall, when firelight shines far more vividly in the dark of night than in daylight---echoing the rabbinic proverb of ``a lamp at midday'' (\latin{sheraga be-\d{t}ihara}). Abarbanel answers the Ran with sharp precision: that explanation holds for \emph{verse 17} (``and when the sun had set and thick darkness came, behold a smoking furnace and a torch of fire''), where the fiery torch actually passes between the pieces, yet it explains nothing in \emph{verse 12}, where no fiery torch is mentioned at all---only the spoken oracle of four hundred years of sojourn and servitude!

The Warsaw 1862 excerpt preserves Abarbanel's opening formula and Questions 7 and 16 without the body of his commentary answering them, so his own later resolution is not extrapolated beyond these two questions. Even from the prefatory questions alone, however, two historical findings stand firm. First, the tension between the stars of Genesis 15:5 and the sunset of Genesis 15:12 was recognized and stated in classic Jewish exegesis more than five centuries ago, directly beside the earlier visionary and anti-astrological traditions Abarbanel cites. Second, the very scholar who stated the waking dilemma in its sharpest form treated waking daytime star-counting not as an exegetical insight, but as an impossibility that ruled out a daytime literal reading of verse 5.
